% --- VI. Theme Extension ---
% Requisitos del enunciado para el tema:
%   - se activa con una opción de línea de comandos; sin ella, todo igual
%   - al menos una prueba con resultado calculado a mano
%   - un experimento (cuenta como uno de los tres del reporte)
\section{Theme: \todo{theme name}}
\label{sec:theme}
After the core simulator was complete, each team extended it with one extra mechanism, called the \emph{theme}, and used an experiment to study what this mechanism really changes~\cite{brief}. The extension is turned on with a command-line option, so a run without the option behaves exactly as before. This keeps all the results of Sections~\ref{sec:validation} and~\ref{sec:experiments} valid.

\subsection{Motivation}
\todo{What problem does the mechanism solve? One paragraph, with a reference to the textbook.}

\subsection{Design}
\todo{What was added to the simulator, in which files, and the command-line option that turns it on. A run without the option behaves exactly as before.}

\begin{lstlisting}
$ ./ossim workload.txt ... --TODO-option
\end{lstlisting}

\subsection{Validation}
\todo{The hand-computed test: the input, the expected result worked out by hand, and the result of the simulator.}

\subsection{Experiment 3: \todo{title}}

\textbf{Question.} \todo{The question the theme experiment answers.}

\textbf{Method.} \todo{Workloads, parameters, what is varied.}

\textbf{Results.} \todo{Describe the plot or table.}

\begin{figure}[H]
    \centering
    \figplaceholder{Plot or table of the theme experiment.}
    \caption{\todo{caption}}
    \label{fig:exp3}
\end{figure}

\textbf{Discussion.} \todo{What the mechanism actually changes, and when it does not help.}
